% Generated draft: requires mathematical and source review.
\documentclass[11pt]{article}
\usepackage[margin=25mm]{geometry}
\usepackage{fontspec}
\setmainfont{Latin Modern Roman}
\usepackage{amsmath,amssymb,mathtools,unicode-math}
\setmathfont{Latin Modern Math}
\usepackage{enumitem,needspace,xurl,hyperref}
\setlength{\parindent}{0pt}
\setlength{\parskip}{5pt}
\setlength{\emergencystretch}{3em}
\newcommand{\N}{\mathbb{N}}
\newcommand{\Z}{\mathbb{Z}}
\newcommand{\Q}{\mathbb{Q}}
\newcommand{\R}{\mathbb{R}}
\newcommand{\C}{\mathbb{C}}
\newcommand{\F}{\mathbb{F}}
\newcommand{\A}{\mathbb{A}}
\newcommand{\PP}{\mathbb{P}}
\newcommand{\E}{\mathbb{E}}
\newcommand{\eps}{\varepsilon}
\newcommand{\dd}{\,\mathrm{d}}
\begin{document}
\section*{UnsolvedMath ID: 30001259}
\textbf{Source problem code:} catalogue entry
\subsection{Definitions and mathematical statement}
For notation, if \(n\geq 1\) and \(\Omega\subseteq\R^n\) is open, then
\[
BV(\Omega)=\{u\in L^1(\Omega):Du\text{ is a finite }\R^n\text{-valued Radon measure}\},
\]
where \(Du\) is the distributional derivative. Its total variation is
\[
|Du|(\Omega)=\sup\left\{\int_\Omega u\,\operatorname{div}\varphi\,\mathrm{d}x:
\varphi\in C_c^1(\Omega;\R^n),\ |\varphi(x)|\leq 1\text{ for every }x\in\Omega\right\}.
\]
These definitions explain the notation; the supplied record does not specify the dimension or hypotheses on the intended domain \(\Omega\).

\textbf{Problem as supplied.} For the variational eigenproblem of the nonsmooth \(1\)-Laplace operator on \(BV(\Omega)\), define and characterize higher eigensolutions beyond the first minimizers, accounting for the ambiguity between a single eigenvalue equation and multiple eigenvalue equations.

\textbf{Unresolved formulation.} The record supplies neither the variational functional and normalization constraint nor the boundary conditions, the notion of solution, or the competing single and multiple eigenvalue equations. Consequently, it does not determine a specific eigenproblem, the meaning of the first minimizers, or an indexing and multiplicity convention for higher eigensolutions. The total variation defined above alone does not specify these choices. A precise quantified characterization therefore requires additional source information; no particular completion of these missing hypotheses is asserted here.

The supplied literature assessment reports that a variational weak-slope construction gives a sequence of higher \(1\)-Laplacian eigensolutions, while the broader question of intrinsic characterization remains unresolved. This is supplied dated curation, not an independently verified mathematical or literature conclusion.
\subsection{Short English statement}
The problem asks how to define and intrinsically characterize higher eigensolutions of the nonsmooth \(1\)-Laplacian beyond the first minimizers, including how single and multiple eigenvalue equations should be interpreted. The supplied record leaves the domain assumptions, boundary conditions, normalization, and equations unspecified, so it does not determine a unique precise eigenproblem. Its supplied status is partially solved.
\subsection{Sources}
\begin{itemize}
\item Local UnsolvedMath problem registry, permanent ID \(30001259\), source problem code \texttt{catalogue entry}, title ``Higher Eigensolutions of the \(1\)-Laplacian.'' The record belongs to the Oberwolfach Reports Open Problems collection and gives the proposed year as 2009; no individual proposer is identified.
\item Supplied report citation: \emph{Topological and Variational Methods for Partial Differential Equations} (2009), Oberwolfach Reports, pp.\ 1427--1486, DOI \texttt{10.4171/owr/2009/26}. \url{https://doi.org/10.4171/owr/2009/26}. The record describes these reports as documenting mathematical workshops and their open-problem discussions; it supplies no problem-specific page or verbatim report passage.
\item Supplied literature source, labeled ``Milbers--Schuricht, sequence of eigensolutions'': \url{https://academic.oup.com/jlms/article/82/1/74/2194080}.
\item Supplied literature source, labeled ``later 1-Laplacian developments'': \url{https://ems.press/content/serial-article-files/46439?nt=1}.
\item Supplied literature source, labeled ``original Oberwolfach report'': \url{https://doi.org/10.4171/owr/2009/26}.
\item Literature assessment date: \texttt{2026-08-21}. The supplied dated curation labels the problem partially solved and reports a weak-slope construction of a sequence of higher eigensolutions, with the broader intrinsic-characterization question unresolved. The background explicitly requests independent verification of time-sensitive status and cited claims. Neither these sources nor this assessment have been independently verified for this draft.
\item The registry labels the statement \texttt{corrected\_verified} and states that the complete statement was preserved while publication formatting was normalized. These are supplied editorial assertions, not independent verification of the formulation.
\end{itemize}
\end{document}
